\chapter{Desarrollo}
\label{cap:desarrollo}

Este capítulo describe el análisis de requisitos, la arquitectura propuesta, las decisiones
de diseño adoptadas y la implementación del sistema. La exposición sigue el mismo orden
incremental en que se ha llevado a cabo el desarrollo, de modo que cada apartado se apoya en
las capacidades construidas en el anterior.

\section{Análisis de requisitos}

\subsection{Requisitos funcionales}

\begin{description}
    \item[RF1] El sistema debe ingerir telemetría de sistema operativo y de red procedente de
    fuentes heterogéneas, sobre un flujo continuo de datos.
    \item[RF2] El sistema debe normalizar los eventos ingeridos a un esquema canónico común,
    preservando la información de origen que no encuentre correspondencia en el modelo.
    \item[RF3] El sistema debe detectar técnicas adversarias individuales sobre el flujo
    normalizado, etiquetando cada detección con la técnica de MITRE ATT\&CK correspondiente.
    \item[RF4] El sistema debe detectar secuencias de ataque multi-etapa, esto es, patrones
    formados por varias detecciones relacionadas temporalmente.
    \item[RF5] El sistema debe mantener un grafo que integre activos, topología,
    vulnerabilidades e inteligencia de amenazas.
    \item[RF6] El sistema debe derivar rutas de ataque factibles hacia los activos críticos
    mediante razonamiento lógico sobre dicho grafo.
    \item[RF7] El sistema debe asignar a cada incidente una puntuación de riesgo acompañada
    de la descomposición de los factores que la determinan.
    \item[RF8] El sistema debe permitir la consulta analítica de los eventos almacenados con
    fines de investigación y evaluación.
    \item[RF9] El sistema debe mantener la trazabilidad de cada evento respecto de su origen,
    de forma que sea posible contrastar los resultados con la verdad conocida durante la
    evaluación.
\end{description}

\subsection{Requisitos no funcionales}

\begin{description}
    \item[RNF1] \textbf{Reproducibilidad}: el sistema completo debe poder desplegarse desde
    cero sobre una máquina de trabajo, sin intervención manual de configuración.
    \item[RNF2] \textbf{Corrección temporal}: el procesamiento debe apoyarse en el tiempo de
    ocurrencia real de los eventos, no en el de su procesamiento.
    \item[RNF3] \textbf{Ausencia de pérdida silenciosa}: todo evento que no pueda procesarse
    debe derivarse a una cola de rechazos con el motivo del fallo, nunca descartarse.
    \item[RNF4] \textbf{Observabilidad}: debe ser posible medir el estado y el rendimiento de
    cada componente del sistema en tiempo real.
    \item[RNF5] \textbf{Contención de recursos}: el sistema debe operar dentro de los límites
    de una estación de trabajo, sin requerir infraestructura dedicada.
    \item[RNF6] \textbf{Trazabilidad de decisiones}: las salidas del sistema deben ser
    explicables, evitando puntuaciones cuyo fundamento no pueda inspeccionarse.
\end{description}

\subsection{Restricciones}

El proyecto se apoya exclusivamente en herramientas de código abierto y conjuntos de datos
públicos. Toda la emulación de comportamiento adversario se realiza sobre datos previamente
capturados en entornos de laboratorio, sin ejecución sobre sistemas en producción. Los
recursos de cómputo se limitan a una estación de trabajo, lo que descarta por diseño
cualquier planteamiento que requiera el entrenamiento de modelos de gran tamaño.

\section{Visión general de la arquitectura}

La arquitectura se organiza en \textbf{dos planos de datos} que convergen en una
\textbf{capa de correlación y decisión}.

El \textbf{plano de telemetría} recorre el camino que va desde los sensores hasta las
detecciones: los eventos crudos se publican en un bus de eventos, se normalizan a un esquema
canónico, se evalúan contra reglas de detección y se correlacionan en busca de secuencias
multi-etapa. El \textbf{plano de inteligencia} construye en paralelo un grafo de conocimiento
a partir de fuentes de CTI y de la base de MITRE ATT\&CK. La capa de correlación cruza ambos:
enriquece las detecciones con el contexto del grafo, razona sobre las rutas de ataque
resultantes y emite la priorización final.

Sobre esa estructura se articulan los tres pilares que constituyen la aportación del trabajo:

\begin{description}
    \item[Pilar A --- Detección multi-etapa] Reconocimiento de secuencias de la
    \textit{kill chain} mediante procesamiento de eventos complejos, en lugar de alertas
    aisladas.
    \item[Pilar B --- Grafo de ataque razonado] Derivación de rutas de ataque factibles hacia
    los activos críticos mediante razonamiento lógico, con trazabilidad de cada paso.
    \item[Pilar C --- Priorización explicable] Asignación a cada incidente de una puntuación
    descompuesta en factores interpretables, acompañada de una justificación legible.
\end{description}

% TODO: insertar figura del diagrama de arquitectura completo (img/arquitectura.pdf)

\subsection{Contrato de flujos de datos}

La comunicación entre etapas se realiza mediante \textit{topics} de un bus de eventos, con
una convención de nombres que refleja el grado de elaboración del dato: el prefijo
\texttt{raw} designa telemetría cruda, \texttt{norm} eventos ya normalizados, \texttt{det}
detecciones y \texttt{enr} incidentes enriquecidos. Los rechazos de cada etapa se derivan a
un \textit{topic} con prefijo \texttt{dlq}.

La Tabla~\ref{tab:topics} recoge el contrato establecido.

\begin{table}[H]
\centering
\begin{tabular}{|l|c|l|l|}
\hline
\textbf{Topic} & \textbf{Part.} & \textbf{Clave} & \textbf{Contenido} \\
\hline
\texttt{raw.host.windows} & 3 & \textit{hostname} & Telemetría cruda de sistema \\
\texttt{raw.net.zeek} & 3 & IP origen & Telemetría cruda de red \\
\texttt{norm.events} & 6 & \textit{hostname} & Eventos normalizados a OCSF \\
\texttt{det.alerts} & 3 & \textit{hostname} & Detecciones de reglas individuales \\
\texttt{det.campaigns} & 3 & \textit{hostname} & Secuencias multi-etapa detectadas \\
\texttt{enr.incidents} & 3 & \textit{id} incidente & Incidentes priorizados \\
\texttt{dlq.*} & 1 & --- & Rechazos, con el motivo del fallo \\
\hline
\end{tabular}
\caption{Contrato de \textit{topics} del sistema.}
\label{tab:topics}
\end{table}

La elección del \textbf{nombre de host como clave de particionado} responde a un requisito
funcional concreto: garantizar que todos los eventos de una misma máquina se procesen en
orden. La detección de secuencias carece de sentido si los eventos que la componen pueden
llegar desordenados, por lo que esta decisión, aparentemente menor, condiciona la viabilidad
del Pilar~A. Como se documenta en la Sección~\ref{sec:incidencia_watermark}, tiene además una
consecuencia no evidente sobre el procesamiento temporal.

\section{Decisiones de diseño}

Esta sección recoge las decisiones de arquitectura relevantes junto con su justificación y
las alternativas descartadas. Se documentan de forma explícita porque constituyen, en buena
medida, el contenido técnico del trabajo: en un proyecto de integración como este, el valor
reside tanto en las piezas elegidas como en los motivos de la elección.

\subsection{Representación del grafo: grafo de propiedades frente a RDF}

El planteamiento inicial contemplaba representar el grafo de conocimiento mediante RDF y
consultarlo con SPARQL, siguiendo la tradición de la web semántica. Esta opción se
\textbf{descartó} en favor de un grafo de propiedades consultado mediante Cypher, por tres
motivos. Primero, el grafo de propiedades es el modelo predominante en las herramientas de
análisis de rutas de ataque del ámbito de la seguridad ofensiva, lo que facilita la
comparación con el estado de la práctica. Segundo, la disponibilidad de bibliotecas de
analítica de grafos (centralidad, caminos mínimos, detección de comunidades) sobre este
modelo es considerablemente mayor. Y tercero, la capacidad de inferencia que aportaría una
ontología OWL queda cubierta —y superada para este dominio— por el razonador Datalog descrito
a continuación, que además es específico de seguridad.

\subsection{Razonamiento: Datalog frente a reglas ontológicas}

Para la derivación de rutas de ataque se adopta un motor Datalog siguiendo el enfoque de
MulVAL \cite{mulval}, en lugar de reglas de inferencia sobre una ontología. La razón es que
Datalog aplicado a análisis de seguridad de red es un enfoque canónico y consolidado, que
proporciona de forma natural la trazabilidad exigida por el requisito RNF6: cada conclusión
derivada puede justificarse indicando la regla aplicada y los hechos que la sustentan.

\subsection{Motor de procesamiento: Flink frente a micro-lotes}

Se adopta Apache Flink como motor de procesamiento en lugar de una solución basada en
micro-lotes. El motivo es que el Pilar~A requiere procesamiento con estado sobre tiempo de
evento y ventanas temporales precisas, capacidades que constituyen el núcleo del modelo de
Flink. Una arquitectura de micro-lotes obligaría a emular ese comportamiento con una
complejidad añadida considerable y garantías más débiles.

Dentro de Flink existían a su vez dos vías para expresar los patrones de secuencia: la
\textit{Pattern API} de FlinkCEP, disponible únicamente desde Java, y la construcción
\texttt{MATCH\_RECOGNIZE} de SQL, accesible también desde Python. Dado que el resto del
sistema se desarrolla en Python y que la segunda opción es un estándar SQL plenamente
expresivo para el tipo de patrones requeridos, se optó por esta última. La decisión se
validó experimentalmente antes de comprometerla, como se describe en la
Sección~\ref{sec:validacion_cep}.

\subsection{Componentes de infraestructura}

El resto de componentes se seleccionaron atendiendo a criterios de eficiencia operativa
sobre una estación de trabajo:

\begin{itemize}
    \item \textbf{Bus de eventos}: se emplea Redpanda, compatible con la API de Kafka pero
    implementado como binario único sin dependencia de JVM ni de un servicio de coordinación
    externo, lo que reduce sensiblemente el consumo de recursos y la complejidad operativa.
    \item \textbf{Almacén analítico}: se emplea ClickHouse, base de datos columnar orientada
    a analítica, por su rendimiento en consultas de agregación y filtrado sobre grandes
    volúmenes de eventos, que es el patrón de acceso característico de la investigación de
    incidentes.
    \item \textbf{Contenedores frente a máquinas virtuales}: todos los servicios se despliegan
    como contenedores sobre un único anfitrión. Dedicar una máquina virtual completa a cada
    servicio habría supuesto multiplicar el consumo de memoria por el sistema operativo
    invitado sin aportar aislamiento adicional relevante para este caso de uso.
\end{itemize}

\subsection{Papel del aprendizaje automático}

El aprendizaje automático se incorpora \textbf{únicamente allí donde aporta capacidades que
el enfoque simbólico no cubre}. Este criterio descarta expresamente los modelos que
duplicarían funciones ya resueltas: un clasificador que replicase el comportamiento de una
regla de detección añadiría opacidad sin ganar cobertura.

Bajo ese principio se identifican cuatro puntos de aplicación:

\begin{enumerate}
    \item \textbf{Detección de anomalías} sobre el flujo normalizado, que complementa a las
    reglas señalando comportamiento inusual sin regla asociada.
    \item \textbf{Extracción de técnicas} a partir de informes de amenazas en lenguaje
    natural, tarea inabordable mediante reglas.
    \item \textbf{Aprendizaje sobre el grafo de ataque}, para sugerir rutas de progresión no
    contempladas explícitamente por las reglas de interacción.
    \item \textbf{Ajuste del modelo de puntuación} a partir de datos, en lugar de pesos
    fijados manualmente.
\end{enumerate}

La restricción operativa es que el entrenamiento debe ser viable sobre la estación de
trabajo descrita, lo que acota el tamaño de los modelos empleados pero no impide el ajuste
fino de modelos del dominio ni el entrenamiento de modelos clásicos y de grafos.

En todos los casos se preserva la explicabilidad mediante el principio expuesto en la
Sección~\ref{sec:ml}: el componente estadístico propone y el simbólico verifica. Ninguna
salida del sistema se presenta al analista sin una justificación inspeccionable.

\section{Entorno de desarrollo}

El desarrollo se realiza sobre una estación de trabajo con procesador de 16 núcleos, 64~GB de
memoria y GPU dedicada, ejecutando el subsistema de Windows para Linux (WSL2) con Ubuntu como
sistema anfitrión de contenedores. Los límites de recursos asignados al subsistema
(44~GB de memoria y 24 hilos) se fijan de forma declarativa, reservando el resto para el
sistema operativo anfitrión.

Los servicios se orquestan mediante Docker Compose con \textbf{perfiles por fase}, lo que
permite levantar únicamente el subconjunto de componentes necesario para la tarea en curso.
Esta organización responde al requisito RNF5: el sistema completo excede los recursos
razonables de una estación de trabajo si todos los servicios se ejecutan simultáneamente,
mientras que por perfiles el consumo se mantiene contenido.

\begin{table}[H]
\centering
\begin{tabular}{|l|l|l|}
\hline
\textbf{Perfil} & \textbf{Servicios} & \textbf{Uso} \\
\hline
\texttt{core} & Bus de eventos y consola & Siempre activo \\
\texttt{stream} & Motor de procesamiento & Desarrollo de detección \\
\texttt{store} & Almacén analítico & Consulta y evaluación \\
\texttt{graph} & Base de datos de grafos & Trabajo sobre el Pilar B \\
\texttt{cti} & Plataforma de inteligencia & Ingesta de CTI \\
\texttt{observability} & Métricas y visualización & Monitorización \\
\hline
\end{tabular}
\caption{Perfiles de despliegue definidos.}
\label{tab:perfiles}
\end{table}

\subsection{Incidencias del entorno}

La construcción del entorno presentó una cadena de incompatibilidades derivadas de emplear
versiones muy recientes del sistema operativo y del intérprete de Python, que se documentan
por su valor práctico:

\begin{itemize}
    \item La distribución de Linux empleada es posterior a las versiones para las que existen
    repositorios oficiales de algunos componentes, lo que obligó a fijar explícitamente el
    repositorio de la última versión con soporte a largo plazo disponible.
    \item La versión de Python del sistema (3.14) carece de paquetes precompilados para las
    bibliotecas de integración con el motor de procesamiento, cuyas dependencias no soportan
    aún esa versión. Se resolvió mediante un entorno virtual independiente con Python 3.12,
    manteniendo el entorno principal para el resto del código.
    \item La biblioteca de integración requiere una máquina virtual Java en el anfitrión, no
    solo en los contenedores, lo que no resulta evidente hasta que falla la primera
    invocación.
\end{itemize}

\section{Ingesta de datos}

\subsection{Selección de los conjuntos de datos}

La validación requiere telemetría de ataques reales con \textbf{verdad conocida}, esto es,
con etiquetado explícito de las técnicas presentes. Se emplean para ello los conjuntos de
datos abiertos de OTRF \cite{otrf}, que reúnen capturas de laboratorio de técnicas
adversarias con dos propiedades especialmente valiosas: cada conjunto está mapeado a las
técnicas de MITRE ATT\&CK correspondientes, y las capturas incluyen no solo los eventos del
ataque sino también la actividad benigna concurrente, lo que permite evaluar falsos
positivos sobre datos realistas.

El repositorio organiza los conjuntos en dos familias. Los \textbf{atómicos} corresponden a
una única técnica aislada y resultan adecuados para la validación de componentes; los
\textbf{compuestos} recogen cadenas de ataque completas de varias etapas, y constituyen el
material de evaluación del Pilar~A.

Para el desarrollo inicial se seleccionó el conjunto atómico correspondiente al volcado de
memoria del proceso LSASS mediante la biblioteca \texttt{comsvcs.dll}, identificado como
\texttt{SDWIN-201018195009} y etiquetado con la técnica \texttt{T1003.001} de la táctica
\texttt{TA0006}. La elección responde a que, pese a tratarse de un conjunto pequeño
(184~eventos), contiene una cadena de acciones representativa y dos fuentes de auditoría
independientes que registran el mismo comportamiento, lo que permite validar la correlación
cruzada.

\subsection{Reproducción de la telemetría}

La ingesta se implementa mediante un componente de reproducción que lee los conjuntos de
datos —almacenados en formato JSON Lines, un evento por línea— y los publica en el bus de
eventos. Su diseño incorpora cuatro decisiones relevantes:

\begin{enumerate}
    \item \textbf{Clave de particionado por host}, conforme al contrato de la
    Tabla~\ref{tab:topics}, para garantizar el orden por máquina.
    \item \textbf{Preservación del tiempo de evento original}: la marca temporal del evento
    se propaga como marca temporal del mensaje en el bus, lo que permite que las etapas
    posteriores operen sobre tiempo de evento y no sobre tiempo de procesamiento (RNF2).
    \item \textbf{Trazabilidad mediante cabeceras}: cada mensaje incorpora el identificador
    del conjunto de datos de origen, lo que satisface el requisito RF9 y permitirá contrastar
    las detecciones con la verdad conocida durante la evaluación.
    \item \textbf{Control de velocidad de reproducción}: el componente admite reproducir a
    máxima velocidad, útil para pruebas, o respetando el ritmo temporal original de la
    captura, útil para demostraciones realistas del comportamiento en tiempo real.
\end{enumerate}

La reproducción del conjunto seleccionado se completó con la totalidad de los 184~eventos
publicados sin pérdidas, a un rendimiento aproximado de 3.300~eventos por segundo.

\section{Normalización a OCSF}

\subsection{Selección del subconjunto del esquema}

OCSF define un número elevado de clases de evento, de las cuales solo un subconjunto resulta
pertinente para la telemetría considerada. La selección se realizó a partir de las clases
efectivamente presentes en los conjuntos de datos, y se organizó en tres oleadas de
implementación para permitir que el sistema fuese funcional desde la primera.

\begin{table}[H]
\centering
\begin{tabular}{|l|c|l|}
\hline
\textbf{Clase OCSF} & \textbf{\texttt{class\_uid}} & \textbf{Fuentes que absorbe} \\
\hline
File System Activity & 1001 & Creación y escritura de ficheros \\
Module Activity & 1005 & Carga de bibliotecas dinámicas \\
Scheduled Job Activity & 1006 & Tareas programadas y servicios \\
Process Activity & 1007 & Creación, terminación y acceso a procesos \\
Event Log Activity & 1008 & Manipulación del registro de auditoría \\
Script Activity & 1009 & Ejecución de guiones \\
Account Change & 3001 & Alta y modificación de cuentas \\
Authentication & 3002 & Inicios de sesión \\
Authorize Session & 3003 & Asignación de privilegios especiales \\
Group Management & 3006 & Modificación de pertenencia a grupos \\
Network Activity & 4001 & Conexiones de red \\
DNS Activity & 4003 & Resolución de nombres \\
\hline
\end{tabular}
\caption{Subconjunto de clases OCSF adoptado.}
\label{tab:ocsf}
\end{table}

Merece una mención específica el tratamiento de los eventos de \textbf{registro de Windows}.
El esquema base de OCSF no define una clase específica para esta categoría, por lo que se
optó por asignarlos a la clase más próxima conservando la totalidad del detalle original en
el campo \texttt{unmapped}. La decisión se documenta como tal, dado que una futura extensión
del esquema podría cubrir este caso de forma más precisa.

\subsection{Arquitectura de los mapeadores}

La normalización se estructura en tres niveles. Un módulo común construye el esqueleto OCSF
compartido por todos los eventos: identificadores de clase y actividad, marca temporal,
metadatos, dispositivo de origen y los dos campos de extensión propios descritos más
adelante. Sobre ese esqueleto, un conjunto de \textbf{mapeadores específicos} —uno por
familia de origen y tipo de evento— añade los campos particulares de cada clase. Finalmente,
un \textbf{despachador} enruta cada evento entrante a su mapeador correspondiente en función
del par formado por el canal de origen y el identificador de evento.

Cuando el despachador no encuentra un mapeador para un evento, lanza una excepción específica
que la etapa de procesamiento traduce en la derivación del evento a la cola de rechazos,
junto con el motivo. Este mecanismo materializa el requisito RNF3: no existe descarte
silencioso.

\subsection{Extensiones propias del esquema}

El modelo incorpora dos extensiones sobre el esquema estándar:

\begin{itemize}
    \item El campo \texttt{unmapped}, previsto por el propio OCSF, conserva toda la
    información de origen sin correspondencia en el modelo canónico.
    \item Un bloque de extensión propio recoge la \textbf{trazabilidad} (identificador del
    conjunto de datos, canal y tipo de evento de origen) y un conjunto de \textbf{señales de
    detección precalculadas}.
\end{itemize}

La segunda extensión merece justificación, pues constituye una decisión de diseño con
consecuencias importantes aguas abajo. Durante el mapeo, el sistema evalúa un conjunto de
condiciones características de comportamiento adversario y las materializa como campos
booleanos: acceso al proceso que custodia las credenciales, acceso con permisos compatibles
con la lectura de memoria, carga de bibliotecas asociadas al volcado de procesos, escritura de
ficheros de volcado, borrado del registro de auditoría o inicio de sesión remoto.

El beneficio de precalcular estas señales en la normalización es doble. Por un lado, evita que
cada regla de detección deba reinterpretar los campos de origen, simplificando notablemente su
expresión. Por otro, y más relevante, constituye la \textbf{semilla de la explicabilidad}
perseguida en el Pilar~C: cuando el sistema priorice un incidente, podrá indicar exactamente
qué señales concretas contribuyeron a ello.

\subsection{Resultados de cobertura}

La cobertura de normalización —proporción de eventos que encuentran mapeador— se adoptó como
métrica de progreso de esta fase. La Tabla~\ref{tab:cobertura} recoge su evolución sobre el
conjunto de datos de referencia conforme se incorporaban mapeadores.

\begin{table}[H]
\centering
\begin{tabular}{|l|c|c|c|}
\hline
\textbf{Iteración} & \textbf{Normalizados} & \textbf{Rechazados} & \textbf{Cobertura} \\
\hline
Mapeadores iniciales (Sysmon 1/3/5/10/11) & 74 & 110 & 40{,}2\,\% \\
+ Módulos y registro (Sysmon 7/12/13/14) & 148 & 36 & 80{,}4\,\% \\
+ Canal Security (11 tipos de evento) & 184 & 0 & 100{,}0\,\% \\
\hline
\end{tabular}
\caption{Evolución de la cobertura de normalización sobre el conjunto \texttt{SDWIN-201018195009} (184 eventos).}
\label{tab:cobertura}
\end{table}

El proceso ilustra una propiedad deseable del diseño: la propia métrica indica qué construir a
continuación, pues los rechazos se agrupan por el par (canal, tipo de evento) que los provoca.
El resultado final son 12~tipos de evento OCSF distintos generados a partir de un único
conjunto de datos.

La equivalencia entre el flujo por lotes y el flujo en \textit{streaming} se verificó
ejecutando la misma lógica en ambos modos y contrastando los resultados, que coincidieron
exactamente.

\subsection{Correlación cruzada entre fuentes}

Un resultado relevante de esta fase es la confirmación independiente del comportamiento
adversario por parte de dos subsistemas de auditoría distintos. El acceso al proceso LSASS
queda registrado simultáneamente por la telemetría extendida de Sysmon y por la auditoría
nativa de Windows mediante los eventos de manipulación de descriptores de objeto.

Que dos mecanismos independientes registren el mismo hecho tiene un valor operativo directo:
constituye evidencia corroborada frente a un posible falso positivo, y ofrece resiliencia
ante la desactivación de una de las fuentes por parte del adversario.

Adicionalmente, la traza de llamadas preservada en el campo \texttt{unmapped} conserva la
secuencia completa de bibliotecas implicadas en la operación de volcado, lo que constituye la
evidencia forense que confirma la técnica concreta empleada.

\subsection{Incidencia: sincronización temporal entre fuentes}
\label{sec:incidencia_temporal}

Durante la validación se detectó un error de consecuencias potencialmente graves que no
producía fallo alguno. Los dos canales de telemetría considerados exponen campos temporales
distintos: la telemetría extendida incorpora un campo con la marca de tiempo en formato UTC
además del campo estándar, mientras que la auditoría nativa solo dispone de este último. La
implementación inicial priorizaba el campo aparentemente más preciso cuando estaba
disponible.

El resultado fue un desfase de dieciséis horas entre ambos canales que reordenaba
completamente la secuencia reconstruida: la etapa final del ataque —el borrado del registro de
auditoría— aparecía como la primera. El error no generaba ningún rechazo ni advertencia, y
solo se hizo visible al inspeccionar la coherencia lógica de la cadena de eventos.

La corrección consistió en apoyar la normalización en el campo temporal \textbf{común y
coherente} entre fuentes, en lugar del más específico. La lección resulta generalizable: en un
sistema que correlaciona múltiples fuentes, la homogeneidad del criterio temporal prevalece
sobre la precisión nominal de cada fuente considerada por separado.

Tras la corrección, la secuencia reconstruida presenta un entrelazado coherente de ambas
fuentes con resolución de milisegundos, condición necesaria para la detección de secuencias
del Pilar~A.

\section{Persistencia analítica}

El almacenamiento de los eventos normalizados persigue dos objetivos: soportar la
investigación interactiva sobre los datos y proporcionar el sustrato para la evaluación
cuantitativa del sistema.

\subsection{Diseño del esquema}

El diseño adopta una estrategia \textbf{híbrida}. Los campos que participan habitualmente en
filtros y agregaciones se promocionan a columnas dedicadas con tipos específicos, mientras que
el evento OCSF completo se conserva íntegro en una columna adicional. De este modo se obtiene
el rendimiento de una base de datos columnar sobre los accesos frecuentes sin renunciar a la
flexibilidad de disponer del dato completo cuando se requiere un campo no previsto.

Las decisiones de modelado más relevantes son las siguientes:

\begin{itemize}
    \item La \textbf{clave de ordenación} se establece sobre el par (host, tiempo), que
    corresponde al patrón de acceso dominante en la investigación de incidentes: qué ocurrió
    en una máquina concreta durante una ventana temporal determinada.
    \item El \textbf{particionado} se realiza por mes, lo que facilita la gestión y el purgado
    del almacén por rangos temporales.
    \item Las \textbf{señales de detección} precalculadas durante la normalización se
    materializan como columnas booleanas, lo que reduce las consultas de investigación sobre
    comportamiento adversario a filtros triviales.
    \item Se aplican \textbf{codificaciones específicas} por columna: diferencial más
    compresión para las marcas temporales, que presentan valores próximos entre sí, y
    diccionario para los campos de baja cardinalidad como nombres de host o de clase.
\end{itemize}

\subsection{Consultas de investigación}

La materialización de las señales como columnas permite reconstruir la cadena completa del
ataque mediante una única consulta declarativa que proyecta cada evento sobre la etapa
correspondiente y los ordena temporalmente. El resultado sobre el conjunto de referencia
recupera la secuencia esperada: preparación mediante carga de bibliotecas de volcado, acceso
al proceso de credenciales confirmado por ambas fuentes de auditoría, escritura del volcado a
disco y borrado del registro de auditoría.

Conviene señalar que esta reconstrucción se obtiene \textbf{antes de aplicar regla de
detección alguna}, únicamente a partir de las señales calculadas durante la normalización, lo
que valida el enfoque adoptado en esa fase.

\subsection{Incidencia: exposición del servicio de almacenamiento}

Se documenta una segunda incidencia por su carácter no evidente. El servicio de
almacenamiento, ejecutado en contenedor con el puerto correspondiente publicado, rechazaba
las conexiones procedentes del anfitrión pese a responder correctamente a las peticiones
originadas dentro del propio contenedor. El síntoma resultaba engañoso, ya que el error
devuelto sugería un problema de red o de redirección de puertos.

El diagnóstico se alcanzó comprobando que el servicio tampoco respondía a peticiones dirigidas
a su propia dirección de red: escuchaba exclusivamente en la interfaz de bucle local interna.
La solución consistió en declarar explícitamente la interfaz de escucha en la configuración
del servidor.

\section{Observabilidad}

La instrumentación del sistema responde al requisito RNF4 y persigue dos finalidades:
verificar durante el desarrollo que cada componente se comporta como se espera, y
proporcionar durante la evaluación las medidas de rendimiento que se reportan en el
Capítulo~\ref{cap:resultados}.

La recolección de métricas se realiza mediante un sistema de series temporales que consulta
periódicamente los puntos de exposición de cada componente. Se instrumentaron el bus de
eventos, ambos nodos del motor de procesamiento y el almacén analítico, activando en cada
caso el mecanismo de exposición correspondiente.

Sobre esa base se construyeron tres cuadros de mando, aprovisionados de forma
\textbf{declarativa} —definidos como ficheros de configuración versionados, no configurados
manualmente— para dar cumplimiento al requisito de reproducibilidad RNF1:

\begin{description}
    \item[Salud del \textit{pipeline}] Estado del bus de eventos, volumen por \textit{topic},
    recursos consumidos por el motor de procesamiento y actividad del almacén analítico.
    Verifica el correcto funcionamiento de la infraestructura.
    \item[Panel de seguridad] Distribución de eventos por clase OCSF, señales de detección
    activas, cronología del ataque reconstruida y procesos más frecuentes. Constituye la
    vista orientada al analista.
    \item[Cobertura y calidad de datos] Cobertura de normalización, distribución por canal y
    tipo de evento, conjuntos de datos cargados e indicadores de integridad temporal. Respalda
    las métricas reportadas en la evaluación.
\end{description}

\section{Validación experimental del enfoque de detección de secuencias}
\label{sec:validacion_cep}

Antes de comprometer el desarrollo del Pilar~A con la vía tecnológica seleccionada, se realizó
una validación experimental específica sobre datos reales. El objetivo era confirmar que la
construcción SQL elegida permitía efectivamente detectar secuencias sobre el flujo procedente
del bus de eventos, respetando el tiempo de evento original.

La prueba consistió en definir un patrón mínimo de dos eventos consecutivos particionado por
host, sobre la telemetría ya ingerida. El resultado inicial fue la ausencia total de
coincidencias, lo que motivó el diagnóstico descrito a continuación.

\subsection{Incidencia: bloqueo del avance temporal por particiones inactivas}
\label{sec:incidencia_watermark}

La causa resultó ser una consecuencia no evidente de la decisión de particionado adoptada. Al
emplear el nombre de host como clave y contener el conjunto de datos un único host, la
totalidad de los eventos se concentra en una sola partición, quedando las restantes vacías.

El motor de procesamiento calcula la marca de agua global como el mínimo de las marcas de agua
de todas sus fuentes. Las particiones que no reciben datos nunca emiten marca de agua alguna,
por lo que el mínimo permanece indefinido y \textbf{el tiempo de evento no avanza jamás}. La
consecuencia es que ninguna ventana temporal llega a considerarse cerrada y el motor no emite
resultado alguno, sin producir error que lo delate.

El diagnóstico se alcanzó por descarte, reduciendo el patrón a un único evento —que debería
haber emitido una coincidencia por cada evento de entrada— y comprobando que tampoco producía
salida, lo que descartaba un problema en la expresión del patrón y señalaba al mecanismo
temporal. La solución consiste en configurar un tiempo de inactividad tras el cual una fuente
sin datos se marca como inactiva y deja de bloquear el cálculo.

Resuelta la incidencia, el patrón de dos eventos emitió correctamente las coincidencias
esperadas, encadenadas y con las marcas temporales originales de la captura. Con ello quedó
validada la vía tecnológica y descartada la necesidad de recurrir a la alternativa prevista
como plan de contingencia.

Esta incidencia ilustra una interacción no trivial entre dos decisiones de diseño
independientes: el particionado por host, adoptado para garantizar el orden de eventos, tiene
como efecto colateral la inutilización del mecanismo de avance temporal cuando el número de
hosts es inferior al de particiones.

% ============================================================================
% TODO: las secciones siguientes se completarán conforme avancen las fases
% ============================================================================

\section{Detección basada en reglas}
\label{sec:deteccion-reglas}

Sobre el flujo de eventos normalizados se construye la primera capa de detección, encargada
de identificar técnicas adversarias individuales. Se apoya en el catálogo público de reglas
Sigma descrito en la Sección~\ref{sec:sigma}, lo que permite aprovechar el trabajo de la
comunidad en lugar de redactar reglas propias.

\subsection{Traducción de las reglas al esquema canónico}

La integración presenta un obstáculo no evidente: las reglas del repositorio público están
escritas contra los nombres de campo \textbf{nativos} de Windows y Sysmon, mientras que el
sistema opera sobre eventos ya normalizados a OCSF. Aplicarlas directamente resulta
imposible, pues los campos que referencian no existen en el modelo canónico.

La biblioteca pySigma \cite{pysigma} prevé este escenario mediante el concepto de
\textit{pipeline de procesamiento}: una especificación declarativa que transforma las reglas
para adaptarlas a un modelo de datos concreto. Existe una implementación oficial para OCSF,
que se tomó como punto de partida.

Su análisis reveló, no obstante, \textbf{tres desviaciones} respecto de la semántica del
estándar que la hacían inaplicable al modelo de este trabajo:

\begin{enumerate}
    \item En los eventos de \textit{acceso a proceso}, el proceso origen se mapea al campo
    correspondiente al \textbf{proceso padre}. La correspondencia es incorrecta: en este tipo
    de evento el proceso origen no engendró al destino, sino que ejecuta una acción sobre él.
    OCSF define para ello el concepto de \textit{actor}, que es la correspondencia adecuada.
    \item Los eventos de \textit{carga de módulo} se clasifican como actividad de proceso con
    actividad genérica, cuando el esquema define una clase específica para ello
    (\textit{Module Activity}).
    \item Los accesos a proceso se etiquetan con la actividad genérica «otros» en lugar de la
    correspondiente a la apertura de un descriptor.
\end{enumerate}

La primera desviación produce detecciones incorrectas; las dos últimas impiden que las reglas
lleguen siquiera a evaluarse, dado que la condición sobre el identificador de tipo de evento
nunca se satisface.

Se optó por \textbf{implementar una pipeline propia} en lugar de extender la oficial. El
motivo es técnico: las condiciones que las pipelines añaden a las reglas se acumulan de forma
conjuntiva, por lo que una corrección superpuesta habría exigido simultáneamente dos valores
distintos del mismo campo, anulando toda posibilidad de coincidencia.

La pipeline resultante cubre nueve categorías de origen y establece la correspondencia entre
los campos nativos y el modelo OCSF del proyecto, dirigiendo al bloque de campos no mapeados
aquellos sin equivalente en el esquema. Su desarrollo constituye, en la práctica, la
implementación de un \textit{backend} de Sigma específico para este sistema.

\subsection{Evaluación de las reglas}

Para aplicar las reglas se plantearon dos vías. La primera consiste en \textbf{generar
consultas} en el lenguaje del almacén analítico y ejecutarlas periódicamente; es directa,
pero opera necesariamente en modo diferido. La segunda consiste en \textbf{evaluar la
estructura de la regla directamente contra cada evento} en memoria, sin lenguaje intermedio,
lo que permite el funcionamiento en tiempo real.

Se implementaron ambas: la primera como validación cruzada y la segunda como componente
operativo del sistema.

\subsubsection{Alcance de los modificadores soportados}

El formato Sigma admite un conjunto amplio de modificadores que alteran la comparación de
valores. Implementarlos todos constituiría un proyecto en sí mismo, por lo que se acotó el
alcance a partir de una medición sobre el repositorio completo de reglas.

\begin{table}[H]
\centering
\begin{tabular}{|l|r|r|}
\hline
\textbf{Modificador} & \textbf{Apariciones} & \textbf{Acumulado} \\
\hline
\texttt{contains}    & 3.362 & 45{,}1\,\% \\
\texttt{endswith}    & 2.494 & 78{,}5\,\% \\
\texttt{all}         &   854 & 90{,}0\,\% \\
\texttt{startswith}  &   431 & 95{,}8\,\% \\
\texttt{re}          &   155 & 97{,}9\,\% \\
\hline
Resto (\texttt{windash}, \texttt{cidr}, otros) & 158 & 100\,\% \\
\hline
\end{tabular}
\caption{Frecuencia de uso de los modificadores de Sigma en el repositorio público.}
\label{tab:modificadores}
\end{table}

Cinco modificadores cubren el 97,9\,\% de los usos reales. El evaluador los implementa
todos, junto con la lógica de composición de condiciones (conjunción, disyunción, negación
y cuantificación). Las construcciones no soportadas se detectan explícitamente y la regla
correspondiente se omite, en lugar de producir una evaluación silenciosamente incorrecta.

\subsubsection{Optimización}

La implementación inicial reconstruía el árbol de condiciones de cada regla en cada
evaluación. Sobre un conjunto de 2.093 reglas, ello suponía repetir más de 385.000 veces un
trabajo que puede realizarse una sola vez por regla.

Se introdujeron dos optimizaciones:

\begin{itemize}
    \item \textbf{Precompilación}: el árbol de condiciones se construye una única vez al
    cargar la regla.
    \item \textbf{Encaminamiento por tipo de evento}: cada regla declara la categoría de
    registro sobre la que aplica, lo que permite descartar de antemano aquellas que no pueden
    coincidir con el tipo del evento evaluado.
\end{itemize}

El efecto combinado se muestra en la Tabla~\ref{tab:optimizacion}. La comprobación esencial
es que el conjunto de detecciones producidas permaneció \textbf{idéntico}: se trata de una
optimización, no de una aproximación.

\begin{table}[H]
\centering
\begin{tabular}{|l|r|r|r|}
\hline
\textbf{Versión} & \textbf{Tiempo} & \textbf{Mejora} & \textbf{Detecciones} \\
\hline
Sin optimizar & 108{,}6 s & --- & 10 \\
Optimizada    & 1{,}49 s & 73$\times$ & 10 \\
\hline
\end{tabular}
\caption{Efecto de las optimizaciones sobre el conjunto de referencia (184 eventos, 2.093 reglas).}
\label{tab:optimizacion}
\end{table}

\subsection{Esquema de alerta}

Cada coincidencia produce una alerta que incorpora la identificación de la regla, su nivel de
severidad, las técnicas y tácticas de ATT\&CK asociadas, el dispositivo y usuario implicados,
la referencia al evento de origen y la trazabilidad al conjunto de datos.

El campo determinante es \texttt{matched\_fields}, que registra \textbf{qué condiciones
concretas dispararon la detección} con sus valores. Constituye la base de la explicabilidad
exigida por el requisito RNF6: la alerta no se limita a afirmar que algo es sospechoso, sino
que aporta la evidencia. A modo de ilustración, sobre el conjunto de referencia el sistema
recupera la línea de comandos completa empleada por el adversario, que coincide literalmente
con la documentada en los metadatos del conjunto de datos.

\subsection{Clasificación de las etiquetas ATT\&CK}

Durante la explotación de los resultados se detectó que las etiquetas del repositorio de
reglas no son homogéneas: bajo un mismo prefijo conviven identificadores de \textbf{técnicas},
de \textbf{software}, de \textbf{grupos de actores} y nombres de \textbf{táctica}, incluyendo
algunos no pertenecientes al estándar y acuñados por la propia comunidad. Una integración que
asuma uniformidad clasifica erróneamente estas etiquetas, contaminando el mapeo ATT\&CK.

El sistema las separa en categorías distintas validando los nombres de táctica contra la
relación oficial, y conserva las no reconocidas en un campo aparte en lugar de descartarlas.

\subsection{Resultados}

\subsubsection{Conjunto de referencia}

Sobre el conjunto atómico de robo de credenciales, el sistema produce \textbf{diez alertas y
ningún falso positivo}, procedentes de cinco reglas independientes pertenecientes a tres
categorías de origen distintas. Todas coinciden en señalar la técnica documentada en los
metadatos del conjunto de datos, lo que confirma la corrección del mapeo.

Resulta relevante que las cinco reglas detectan el mismo ataque por vías diferentes: unas por
el acceso a la memoria del proceso de credenciales, otras por la línea de comandos empleada,
otras por el fichero de volcado escrito a disco. Esta redundancia entre fuentes independientes
constituye evidencia corroborada frente a un posible falso positivo.

\subsubsection{Escenario multi-etapa}

Sobre el escenario de emulación de campaña completa, el sistema procesa 195.197 eventos
normalizados y produce \textbf{741 alertas} (0,38\,\% de los eventos), sin ningún rechazo. Las
alertas implican \textbf{cuatro equipos} distintos y cubren \textbf{33 técnicas} de ATT\&CK,
incluyendo robo de credenciales, persistencia, ejecución mediante servicios y exfiltración.

La distribución temporal es reveladora: aproximadamente el 75\,\% de las alertas se concentra
en el 20\,\% del intervalo analizado, lo que se corresponde con la ventana en que se ejecuta
el grueso de la campaña.

\subsubsection{Fatiga de alertas: evidencia empírica}

El análisis de las alertas producidas ilustra con datos propios el problema que motiva este
trabajo. Una única regla genera 292 alertas, el 39\,\% del total. Su inspección revela que
todas proceden de \textbf{una sola acción del adversario}: el despliegue de un conjunto de
herramientas administrativas mediante un intérprete de comandos, que generó una alerta por
cada uno de los ficheros extraídos, en un intervalo de ocho minutos y sobre un único equipo.

La detección es \textbf{correcta}: el comportamiento observado es efectivamente sospechoso. Lo
inadecuado es la granularidad. Un analista no necesita 292 avisos, sino uno que informe de que
se desplegó un conjunto de herramientas en un equipo determinado.

Este caso concreto cuantifica el problema: una acción del adversario consume 292 unidades de
atención del analista. Constituye asimismo el caso de uso más claro para la correlación
descrita en la Sección siguiente.

\subsubsection{Potencial de reducción de alertas}

Como referencia previa a la implementación de la correlación multi-etapa, se estimó el efecto
de una agrupación elemental por regla, equipo y ventana temporal de diez minutos. El resultado
reduce las 741 alertas a \textbf{87 grupos}, un \textbf{88,3\,\%} de reducción, con un grupo
máximo de 292 elementos.

Esta cifra constituye el \textbf{suelo} contra el que debe medirse el Pilar~A: una heurística
trivial ya alcanza ese nivel, por lo que la correlación basada en patrones debe superarlo
agrupando además alertas de reglas distintas que constituyen una misma cadena de ataque.

\subsubsection{Cobertura ATT\&CK}

Se generó la matriz de cobertura en el formato del ATT\&CK Navigator \cite{navigator}, con dos
capas: la cobertura \textit{potencial} —técnicas para las que existe al menos una regla
cargada— y la \textit{observada} —técnicas efectivamente detectadas sobre los conjuntos de
datos procesados.

\begin{table}[H]
\centering
\begin{tabular}{|l|r|}
\hline
\textbf{Indicador} & \textbf{Valor} \\
\hline
Reglas cargadas & 2.301 \\
Reglas sin etiqueta ATT\&CK & 238 (10,3\,\%) \\
Técnicas con cobertura potencial & 300 \\
\quad de las cuales subtécnicas & 177 \\
Técnicas observadas en los datos & 34 \\
Tasa de activación & 11,3\,\% \\
\hline
\end{tabular}
\caption{Cobertura ATT\&CK del sistema.}
\label{tab:cobertura-attack}
\end{table}

La baja tasa de activación no indica una limitación del sistema, sino de los conjuntos de
datos empleados: existen 266 técnicas con detección disponible que ninguno de los escenarios
procesados llega a ejercitar. Esta observación motiva la construcción de un entorno de
emulación propio que permita activar dicha cobertura latente de forma dirigida.

\subsubsection{Rendimiento}

Sobre el escenario completo, el evaluador procesa 78 eventos por segundo frente a 2.093 reglas
simultáneas, lo que equivale a aproximadamente 163.000 evaluaciones de regla por segundo. El
consumidor está implementado sobre un grupo de Kafka sin asignación fija de particiones, de
modo que el rendimiento es escalable mediante la ejecución de instancias adicionales.

\section{Detección de anomalías}
\label{sec:anomalias}

Las reglas descritas en la sección anterior detectan comportamiento \textbf{documentado}. Para
complementarlas se incorpora un detector de anomalías, cuyo cometido es señalar actividad
inusual aunque no exista regla que la describa.

\subsection{Papel dentro del sistema}

Una decisión de diseño relevante es que la anomalía \textbf{no genera alerta por sí misma}.
La razón es su elevada tasa de falsos positivos: emitir un aviso por cada desviación
estadística agravaría precisamente el problema que el trabajo pretende resolver. Su salida se
emplea en cambio para elevar la confianza de una secuencia ya detectada y para marcar
actividad que merece revisión en ausencia de regla aplicable.

\subsection{Nivel de agregación y características}

El modelo opera sobre \textbf{ventanas de actividad por equipo} en lugar de sobre eventos
individuales, por su mejor correspondencia con la naturaleza multi-etapa del problema: un
ataque se manifiesta como un patrón de actividad, no como un evento aislado.

Se definen 42 características, separadas deliberadamente en dos familias:

\begin{itemize}
    \item \textbf{Absolutas} (26): recuentos de eventos por clase, tipo de actividad y señal,
    así como medidas de diversidad. Dependen del volumen de actividad de la ventana.
    \item \textbf{Relativas} (16): proporciones de cada tipo de actividad sobre el total.
    Son \textbf{independientes del volumen por construcción}.
\end{itemize}

Esta separación no es accesoria: permite determinar experimentalmente qué parte de la
capacidad de detección procede del volumen de actividad y cuál del perfil de comportamiento.

\subsection{Protocolo experimental}

El entrenamiento emplea Isolation Forest \cite{isolationforest} sobre la implementación de
scikit-learn \cite{scikitlearn}, y se comparan dos estrategias:

\begin{enumerate}
    \item \textbf{Entrenamiento con ventanas limpias}: únicamente aquellas sin indicios
    objetivos de actividad maliciosa. Corresponde al planteamiento canónico de aprender el
    comportamiento normal.
    \item \textbf{Entrenamiento con contaminación}: todas las ventanas, declarando la
    proporción esperada de anomalías.
\end{enumerate}

Se incorpora asimismo una variante \textbf{supervisada} mediante bosques aleatorios, con y sin
compensación del desbalanceo por SMOTE \cite{smote}.

El conjunto de evaluación se mantiene \textbf{estrictamente separado} del de entrenamiento: se
reserva un escenario completo, para el que se dispone de las alertas de la capa de reglas como
referencia externa. Todos los experimentos se registran en MLflow \cite{mlflow} con sus
parámetros, métricas y semillas, garantizando su reproducibilidad.

\subsection{Controles experimentales}

Junto a los modelos se incorporan dos controles cuya función es delimitar el mérito real de
los resultados:

\begin{itemize}
    \item \textbf{Control de volumen}: un modelo entrenado con una \textbf{única}
    característica, el número de eventos de la ventana. Si alcanzase un rendimiento comparable
    al del modelo completo, quedaría demostrado que la aportación del enfoque multivariante es
    marginal.
    \item \textbf{Ablación de características potencialmente circulares}: se excluyen aquellas
    que participan en la definición de «ventana limpia», para descartar que el modelo
    aproveche información derivada de su propio criterio de entrenamiento.
\end{itemize}

\subsection{Resultados}

\begin{table}[H]
\centering
\begin{tabular}{|l|c|c|c|}
\hline
\textbf{Configuración} & \textbf{Caract.} & \textbf{AUC-ROC} & \textbf{AUC-PR} \\
\hline
No supervisado, limpias, todas & 42 & \textbf{0{,}995} & 0{,}996 \\
No supervisado, limpias, sin circulares & 34 & 0{,}995 & 0{,}996 \\
\textit{Control: solo volumen} & 1 & \textit{0{,}941} & \textit{0{,}920} \\
Supervisado sin SMOTE, sin circulares & 34 & 0{,}923 & 0{,}882 \\
Supervisado con SMOTE, solo ratios & 16 & 0{,}913 & 0{,}869 \\
No supervisado, contaminado, todas & 42 & 0{,}913 & 0{,}919 \\
No supervisado, limpias, solo ratios & 16 & 0{,}867 & 0{,}806 \\
Supervisado con SMOTE, sin circulares & 34 & 0{,}864 & 0{,}829 \\
Supervisado sin SMOTE, solo ratios & 16 & 0{,}854 & 0{,}798 \\
No supervisado, contaminado, solo ratios & 16 & 0{,}533 & 0{,}646 \\
\hline
\end{tabular}
\caption{Resultados del detector de anomalías sobre el escenario de evaluación.}
\label{tab:anomalias}
\end{table}

De la Tabla~\ref{tab:anomalias} se desprenden cuatro observaciones:

\begin{enumerate}
    \item \textbf{No existe circularidad.} Excluir las características que participan en la
    definición de ventana limpia no degrada el resultado.

    \item \textbf{Volumen y perfil son fuentes de información complementarias.} El control
    basado únicamente en el volumen alcanza 0,941, y el modelo que emplea exclusivamente
    características relativas —independientes del volumen— alcanza 0,867. Ambas discriminan de
    forma independiente, y su combinación (0,995) supera a cualquiera por separado, lo que
    justifica el enfoque multivariante frente a una regla de umbral.

    \item \textbf{La pureza del conjunto de entrenamiento resulta determinante.} Con idéntico
    conjunto de características, entrenar con ventanas limpias produce 0,867 mientras que
    hacerlo con datos contaminados produce 0,533, indistinguible del azar.

    \item \textbf{El efecto de SMOTE depende de la dimensionalidad.} Mejora el resultado con
    16 características (0,854 $\rightarrow$ 0,913) y lo degrada con 34 (0,923 $\rightarrow$
    0,864), comportamiento coherente con la interpolación entre vecinos en espacios de
    dimensión elevada.
\end{enumerate}

\subsection{Limitaciones}

Los resultados deben interpretarse con dos cautelas explícitas.

La primera es el \textbf{tamaño del conjunto de evaluación}: veintiocho ventanas, de las cuales
quince positivas. Con esa magnitud, la incertidumbre asociada a las métricas es
considerable y valores próximos a la unidad deben leerse como indicio de separabilidad
sencilla antes que como evidencia de un modelo excelente.

La segunda, más relevante, afecta a la \textbf{representatividad de los datos}. Las ventanas
que contienen actividad maliciosa presentan una mediana de eventos ocho veces superior a las
benignas, con solapamiento reducido. Esta separabilidad es en buena medida un artefacto de la
naturaleza de los conjuntos de datos empleados: capturas de laboratorio en las que las
técnicas se ejecutan de forma concentrada sin actividad de fondo sostenida. En un entorno
real, la actividad legítima presenta volúmenes comparables y un adversario competente opera
precisamente con discreción.

Ambas limitaciones convergen en la misma necesidad: disponer de telemetría con actividad
benigna prolongada y de volumen comparable al de la actividad maliciosa. Esta constatación
motiva la construcción del entorno de emulación descrito más adelante.

\section{Pilar A: detección de secuencias multi-etapa}
% TODO (F3): modelado de los escenarios a partir de los planes de emulación; definición de los
% patrones de kill chain; tratamiento de las secuencias que cruzan varios hosts mediante
% correlación jerárquica en dos niveles; esquema de campaña emitido.

\section{Plano de inteligencia de amenazas}
% TODO (F4): despliegue de la plataforma CTI e importación de ATT&CK; sincronización de
% fuentes de indicadores; extracción automática de técnicas a partir de informes en lenguaje
% natural y carga del resultado en el grafo.

\section{Pilar B: grafo de ataque y razonamiento}
% TODO (F5): modelo de datos del grafo; poblado del inventario y las vulnerabilidades;
% analítica de grafos para criticidad y caminos; reglas de interacción en Datalog; exportación
% de hechos y reimportación de las rutas derivadas; trazabilidad de la derivación.

\section{Pilar C: priorización explicable y servicios}
% TODO (F6): modelo de puntuación y justificación de los pesos; análisis de sensibilidad;
% generación de la narrativa; interfaz de programación de la aplicación.
